% PAGES: 96-102
% Section 2.9 "References" ends before page 96 (owned by a different
% agent's file); the top of page 96 opens directly with the "2.10
% Exercises" heading, so this file begins fresh with that section rather
% than continuing anything. All sixteen exercises (pages 96-101) are kept
% together in a single \exercises environment so the shared enumerate
% counter is never reset mid-list. Page 102 (the final page of this
% assignment) is a blank transitional verso -- see the note at the very
% end of this file.

\section{Exercises}

\begin{exercises}

\item Let
\[
L_1 \;=\; \begin{pmatrix} 1 & 0 & 0 & 0\\ -1 & 2 & 0 & 0\\ 2 & -2 & 3 & 0\\ 2 & 2 & -3 & 4\end{pmatrix};
\qquad
U_1 \;=\; \begin{pmatrix} 2 & -1 & 0 & 1\\ 0 & -1/2 & 1/2 & 0\\ 0 & 0 & 1/3 & -1/3\\ 0 & 0 & 0 & -1/4\end{pmatrix}.
\]
\[
L_2 \;=\; \begin{pmatrix} 1 & 0 & 0 & 0\\ -1 & 1 & 0 & 0\\ 2 & -1 & 1 & 0\\ 2 & 1 & -1 & 1\end{pmatrix};
\qquad
U_2 \;=\; \begin{pmatrix} 2 & -1 & 0 & 1\\ 0 & -1 & 1 & 0\\ 0 & 0 & 1 & -1\\ 0 & 0 & 0 & -1\end{pmatrix}.
\]
\begin{subparts}
\item Show that $L_1U_1 = L_2U_2$.
\item Explain why this does not contradict the uniqueness of the LU decomposition of a matrix.
\item Show that
\[
L_1 \;=\; L_2\begin{pmatrix} 1&0&0&0\\ 0&2&0&0\\ 0&0&3&0\\ 0&0&0&4\end{pmatrix}
\quad \text{and} \quad
U_1 \;=\; \begin{pmatrix} 1&0&0&0\\ 0&1/2&0&0\\ 0&0&1/3&0\\ 0&0&0&1/4\end{pmatrix}U_2.
\]
\end{subparts}

\item Let $L_1$ and $L_2$ be nonsingular lower triangular matrices and let $U_1$ and $U_2$ be
nonsingular upper triangular matrices. If $L_1U_1 = L_2U_2$, show that there exists
a nonsingular diagonal matrix $D$ such that
\[
L_1 \;=\; L_2D \quad \text{and} \quad U_1 \;=\; D^{-1}U_2.
\]

\item Let
\[
A \;=\; \begin{pmatrix} 2 & -1 & 0 & 1\\ -2 & 0 & 1 & -1\\ 4 & -1 & 0 & 1\\ 4 & -3 & 0 & 2\end{pmatrix}.
\]
\begin{subparts}
\item Show that the LU decomposition with row pivoting of the matrix $A$ is given
by $PA = LU$, where
\[
P \;=\; \begin{pmatrix} 0&0&1&0\\ 0&0&0&1\\ 0&1&0&0\\ 1&0&0&0\end{pmatrix};
\]
\[
L \;=\; \begin{pmatrix} 1&0&0&0\\ 1&1&0&0\\ -0.5&0.25&1&0\\ 0.5&0.25&0&1\end{pmatrix};
\qquad
U \;=\; \begin{pmatrix} 4&-1&0&1\\ 0&-2&0&1\\ 0&0&1&-0.75\\ 0&0&0&0.25\end{pmatrix}.
\]
\item Solve $Ax = b$, where $b = \begin{pmatrix} 3\\ -1\\ 0\\ 2\end{pmatrix}$.
\item Find $A^{-1}$, the inverse matrix of $A$.
\end{subparts}

\item Let $A = \begin{pmatrix} 2&-1&3&-1\\ 1&0&-2&-4\\ 3&1&1&-2\\ -4&1&0&2\end{pmatrix}$ and
$b = \begin{pmatrix} -1\\ 0\\ 1\\ 2\end{pmatrix}$.
\begin{subparts}
\item Find the LU decomposition with row pivoting of the matrix $A$;
\item Use the linear solver \textit{linear\_solve\_lu\_row\_pivoting} to solve the linear
system $Ax = b$.
\item Let $P_1 = \begin{pmatrix} 1&0&0&0\\ 0&0&0&1\\ 0&1&0&0\\ 0&0&1&0\end{pmatrix}$ and
$P_2 = \begin{pmatrix} 0&1&0&0\\ 0&0&0&1\\ 1&0&0&0\\ 0&0&1&0\end{pmatrix}$ be permutation
matrices, and let
\[
L_1 \;=\; \begin{pmatrix} 1&0&0&0\\ 0&1&0&0\\ -0.6667&-0.5833&1&0\\ 0.3333&-0.5833&0.1429&1\end{pmatrix};
\]
\[
U_1 \;=\; \begin{pmatrix} 3&2&-1&-1\\ 0&-4&2&1\\ 0&0&-3.5&-0.0833\\ 0&0&0&1.9286\end{pmatrix}.
\]
Show that
\[
P_1\,A\,P_2 \;=\; L_1U_1.
\]
\item Use forward substitution to solve $L_1y = P_1b$, and use backward substitution to
solve $U_1x_1 = y$. Let $x_2 = P_2x_1$. Show that $x_2$ is the same as the solution
$x$ to $Ax = b$ obtained at (ii), and explain why this happens.
\end{subparts}
\noindent Note: The LU decomposition with full pivoting of $A$ is of the form
$P_1AP_2 = L_1U_1$.

\item Find the LU decomposition without pivoting of the matrix
\[
\begin{pmatrix}
1 & 0 & 0 & 0 & 1\\
-1 & 1 & 0 & 0 & 1\\
-1 & -1 & 1 & 0 & 1\\
-1 & -1 & -1 & 1 & 1\\
-1 & -1 & -1 & -1 & 1
\end{pmatrix}.
\]
\noindent Note: This is the $5\times 5$ version of the classic example of a matrix whose
LU decomposition is unstable.

\item Let
\[
A \;=\; \begin{pmatrix} 2 & -1 & 1\\ -2 & 1 & 3\\ 4 & 0 & -1\end{pmatrix}.
\]
\begin{subparts}
\item Show that the $2\times 2$ leading principal minor of $A$ is $0$, i.e., show that
\[
\det\begin{pmatrix} 2 & -1\\ -2 & 1\end{pmatrix} \;=\; 0.
\]
\item Attempt to do the LU decomposition without pivoting of the matrix $A$, and show
that the division by $U(2,2)$ cannot be performed when trying to compute the second
row of $L$.
\item Show that the matrix $A$ is nonsingular, and compute the LU decomposition with
row pivoting of $A$.
\end{subparts}

\item The LU decomposition with column pivoting of an $n\times n$ nonsingular matrix
$A$ is $AP = LU$, where $P$ is an $n\times n$ permutation matrix, $L$ is an
$n\times n$ lower triangular matrix with all entries on the main diagonal equal to
$1$, and $U$ is an $n\times n$ upper triangular matrix. Write a pseudocode for
solving linear systems of the form $Ax=b$ by using the LU decomposition with column
pivoting of $A$.

\item Write the pseudocode for the forward substitution corresponding to a lower
triangular banded matrix of band $m$, i.e., for solving $Lx=b$ where $b$ is an
$n\times 1$ vector and $L$ is an $n\times n$ lower triangular matrix such that
\[
L(j,k) \;=\; 0, \quad \forall\, 1\le j,k\le n \text{ with } j-k>m.
\]
What is the corresponding operation count?

\item Write the pseudocode for the backward substitution corresponding to an upper
triangular banded matrix of band $m$, i.e., for solving $Ux=b$ where $b$ is an
$n\times 1$ vector and $U$ is an $n\times n$ upper triangular matrix such that
\[
U(j,k) \;=\; 0, \quad \forall\, 1\le j,k\le n \text{ with } k-j>m.
\]
What is the corresponding operation count?

\item Write the pseudocode for the LU decomposition without pivoting for banded
matrices of band $m$. What is the corresponding operation count?

Use the fact that the $L$ and $U$ factors from the LU decomposition without pivoting
of a banded matrix of band $m$ are a banded lower triangular matrix of band $m$ and
a banded upper triangular matrix of band $m$, respectively.

\item What is the operation count for solving a linear system corresponding to a
banded matrix of band $m$ using a linear solver based on the LU decomposition
without pivoting of the matrix?

\item The values of the following coupon bonds with face value \$100 are given:

\begin{center}
\begin{tabular}{ccc}
Bond Type & Coupon Rate & Bond Price \\[0.4em]
10 months semiannual & 3\% & \$101.30 \\
16 months semiannual & 4\% & \$102.95 \\
22 months annual & 6\% & \$107.35 \\
22 months semiannual & 5\% & \$105.45 \\
\end{tabular}
\end{center}

\begin{subparts}
\item List the cash flows and cash flow dates for each bond.
\item Identify the $4\times 4$ matrix and the $4\times 1$ right hand side vector
corresponding to the linear system whose solution are the 4 months, 10 months, 16
months, and 22 months discount factors.
\item Find the 4 months, 10 months, 16 months, and 22 months discount factors.
\end{subparts}

\item The values of the following coupon bonds with face value \$100 are given:

\begin{center}
\begin{tabular}{ccc}
Bond Type & Coupon Rate & Bond Price \\[0.4em]
6 months semiannual & 0 & \$98.50 \\
1 year semiannual & 3\% & \$101.00 \\
18 months semiannual & 5\% & \$102.00 \\
2 years semiannual & 3\% & \$103.50 \\
\end{tabular}
\end{center}

\begin{subparts}
\item List the cash flows and cash flow dates for each bond.
\item Find the 6 months, 1 year, 18 months, and 2 years discount factors.
\end{subparts}

\item The values of the following coupon bonds with face value \$100 are given:

\begin{center}
\begin{tabular}{ccc}
Bond Type & Coupon Rate & Bond Price \\[0.4em]
9 months semiannual & 2\% & \$100.80 \\
15 months semiannual & 4\% & \$103.50 \\
15 months annual & 5\% & \$107.50 \\
21 months semiannual & 5\% & \$110.50 \\
\end{tabular}
\end{center}

\begin{subparts}
\item List the cash flows and cash flow dates for each bond.
\item Find the 3 months, 9 months, 15 months, and 21 months discount factors.
\end{subparts}

\item The following discount factors are obtained by fitting market data:

\begin{center}
\begin{tabular}{cc}
Date & Discount Factor \\[0.4em]
2 months & 99.80 \\
5 months & 99.35 \\
11 months & 98.20 \\
15 months & 97.75 \\
\end{tabular}
\end{center}

The overnight rate is $1\%$.

\begin{subparts}
\item What is the linear system that has to be solved for the cubic spline
interpolation of the zero rate curve?
\item Use cubic spline interpolation to find a zero rate curve for all times less
than 15 months matching the discount factors above.
\item Find the value of a 13 months quarterly bond with $2.5\%$ coupon rate.
\end{subparts}
\noindent Note: A quarterly coupon bond with face value \$100, coupon rate $C$, and
maturity $T$ pays the holder of the bond a coupon payment equal to
$\frac{C}{4}\cdot 100$ every three months, except at maturity. The final payment at
maturity $T$ is equal to the face value of the bond plus one coupon payment, i.e.,
$100 + \frac{C}{4}100$.

\item Consider three assets with the following expected values, standard
deviations, and correlations of their returns:
\begin{align*}
\mu_1 &= 0.10; & \sigma_1 &= 0.15; & \rho_{1,2} &= -0.25; \\
\mu_2 &= 0.15; & \sigma_2 &= 0.30; & \rho_{2,3} &= 0.20; \\
\mu_3 &= 0.20; & \sigma_3 &= 0.35; & \rho_{1,3} &= 0.30.
\end{align*}
\begin{subparts}
\item Find the covariance matrix $M$ of the returns of the three assets.

\noindent Hint: From (7.15), it follows that the covariance matrix is given by
\[
M \;=\; \begin{pmatrix}
\sigma_1^2 & \sigma_1\sigma_2\rho_{1,2} & \sigma_1\sigma_3\rho_{1,3}\\
\sigma_1\sigma_2\rho_{1,2} & \sigma_2^2 & \sigma_2\sigma_3\rho_{2,3}\\
\sigma_1\sigma_3\rho_{1,3} & \sigma_2\sigma_3\rho_{2,3} & \sigma_3^2
\end{pmatrix}.
\]
\item A minimum variance portfolio with $16\%$ expected rate of return and fully
invested in the three assets (i.e., with no cash position) can be set up by
investing a percentage $w_i$ of the total value of the portfolio in asset $i$, with
$i=1:3$, where $w_1$, $w_2$, $w_3$ can be found by solving the following linear
system:
\begin{equation}
\begin{pmatrix} 2M & \bfone & \mu\\ \bfone^t & 0 & 0\\ \mu^t & 0 & 0\end{pmatrix}
\begin{pmatrix} w\\ \lambda_1\\ \lambda_2\end{pmatrix}
\;=\;
\begin{pmatrix} 0\\ 1\\ \mu_P\end{pmatrix},
\end{equation}
where $w = \begin{pmatrix} w_1\\ w_2\\ w_3\end{pmatrix}$ and
\[
\mu_P = 0.16; \quad
\mu = \begin{pmatrix} 0.1\\ 0.15\\ 0.2\end{pmatrix}; \quad
\bfone = \begin{pmatrix} 1\\ 1\\ 1\end{pmatrix}.
\]

Show that the matrices from the LU decomposition with row pivoting of the matrix on
the left hand side of (2.100) are
\[
P \;=\; \begin{pmatrix}
0&0&0&1&0\\
0&1&0&0&0\\
0&0&1&0&0\\
1&0&0&0&0\\
0&0&0&0&1
\end{pmatrix};
\]
\[
L \;=\; \begin{pmatrix}
1 & 0 & 0 & 0 & 0\\
-0.0225 & 1 & 0 & 0 & 0\\
0.0315 & 0.051852 & 1 & 0 & 0\\
0.045 & -0.333333 & 0.038067 & 1 & 0\\
0.1 & 0.246914 & 0.400056 & -0.482738 & 1
\end{pmatrix};
\]
\[
U \;=\; \begin{pmatrix}
1 & 1 & 1 & 0 & 0\\
0 & 0.2025 & 0.0645 & 1 & 0.15\\
0 & 0 & 0.210555 & 0.948148 & 0.192222\\
0 & 0 & 0 & 1.297240 & 0.142683\\
0 & 0 & 0 & 0 & -0.045059
\end{pmatrix}.
\]
\item Find the weights of each asset in the minimum variance portfolio with $16\%$
expected return. Find the standard deviation of the return of this portfolio.
\item Show that the two portfolios below have $16\%$ expected return, and compute
the standard deviation of the returns of each portfolio:
\begin{itemize}
\item $30\%$ invested in asset 1, $20\%$ invested in asset 2, $50\%$ invested in
asset 3;
\item $50\%$ invested in asset 1, $70\%$ invested in asset 3, and short an amount
equal to $20\%$ of the value of the portfolio of asset 2.
\end{itemize}

\noindent Hint: If $w_1$, $w_2$, and $w_3$ denote the weights of asset 1, of asset
2, and of asset 3 in the portfolio, respectively, then the expected return of the
portfolio is
\[
w_1\mu_1 + w_2\mu_2 + w_3\mu_3.
\]
\end{subparts}

\end{exercises}

% PDF page 102 (the last page of this assignment) is blank -- a
% transitional blank verso before Chapter 3 "The Arrow--Debreu one period
% market model," which begins on PDF page 103. Nothing to transcribe on
% that page.
